%% Hand-authored TikZ fragment -- NOT generated by maestro2tex.
%%
%% The electrochemical cell is drawn as an ABSTRACTED symbol everywhere in this
%% chapter: a plain circle with one short stub per electrode and nothing inside.
%% The internal R/C network is not a drawing any more -- the element values in
%% Table~\ref{tab:analog-models} are the definition of the model.
%%
%% Geometry taken read-only from the owner's Virtuoso symbol
%% (castalia/randles_cell, view symbol, 2026-08-22):
%%   ellipse  bBox ((0.2875 -0.5875) (1.0875 0.2125))  -> circle, radius 0.4
%%   stubs    three lines, each from 0.9375r to 1.5625r, terminal name beside it
%%   angles   the two collinear electrodes opposite each other (CE, WE), the
%%            third (RE) perpendicular
%% The stub lengths and the label offsets below are those ratios. The symbol is
%% turned 90 degrees from the OA drawing (which has RE left, CE top, WE bottom)
%% so that CE enters from the left and WE leaves to the right: that is the
%% direction the bench figures run, and it keeps every wire orthogonal. The
%% relative arrangement is unchanged.
%%
%% Interface:
%%   \CellSymThree{<centre>}{<radius>}            CE left, RE top, WE right
%%   \CellSymTwo{<centre>}{<radius>}{<a>}{<b>}    two-terminal, a left, b right
%%   \CellSymTwoV{<centre>}{<radius>}{<a>}{<b>}   two-terminal, a top, b bottom
%% <centre> is a bare coordinate pair (e.g. 6.0,0), <radius> a length in cm, and
%% the label arguments may be empty to draw the stub unlabelled.
%%
%% Each macro leaves coordinates at the stub ENDS for the figure's wiring to
%% meet, plus one at the centre:
%%   (cs-ce) (cs-re) (cs-we) (cs-c)   /   (cs-a) (cs-b) (cs-c)
%% Labels are set in the host figure's own `nlab' style, so the font follows the
%% figure rather than being fixed here; line weight is the figure's default, so
%% the symbol matches the wires it joins.
%%
%% \input this from every figure that draws a cell, after the usual
%% \providecommand{\MaestroRoot}{include/analog/}. The definitions are \gdef so
%% they survive whatever group the figure sits in, and re-inputting the file
%% simply redefines them identically -- no guard needed.

\gdef\CellSymThree#1#2{%
  \pgfmathsetmacro{\csRad}{#2}%
  \pgfmathsetmacro{\csLead}{1.5625*(#2)}%
  \pgfmathsetmacro{\csLbl}{1.28*(#2)}%
  \begin{scope}[shift={(#1)}]
    \draw (0,0) circle[radius=\csRad cm];
    \draw (-\csRad,0) -- (-\csLead,0);
    \draw (0,\csRad) -- (0,\csLead);
    \draw (\csRad,0) -- (\csLead,0);
    \node[nlab, anchor=south] at (-\csLbl,0) {\texttt{CE}};
    \node[nlab, anchor=west]  at (0,\csLbl)  {\texttt{RE}};
    \node[nlab, anchor=south] at (\csLbl,0)  {\texttt{WE}};
    \coordinate (cs-c)  at (0,0);
    \coordinate (cs-ce) at (-\csLead,0);
    \coordinate (cs-re) at (0,\csLead);
    \coordinate (cs-we) at (\csLead,0);
  \end{scope}%
}

\gdef\CellSymTwo#1#2#3#4{%
  \pgfmathsetmacro{\csRad}{#2}%
  \pgfmathsetmacro{\csLead}{1.5625*(#2)}%
  \pgfmathsetmacro{\csLbl}{1.28*(#2)}%
  \begin{scope}[shift={(#1)}]
    \draw (0,0) circle[radius=\csRad cm];
    \draw (-\csRad,0) -- (-\csLead,0);
    \draw (\csRad,0) -- (\csLead,0);
    \node[nlab, anchor=south] at (-\csLbl,0) {#3};
    \node[nlab, anchor=south] at (\csLbl,0)  {#4};
    \coordinate (cs-c) at (0,0);
    \coordinate (cs-a) at (-\csLead,0);
    \coordinate (cs-b) at (\csLead,0);
  \end{scope}%
}

\gdef\CellSymTwoV#1#2#3#4{%
  \pgfmathsetmacro{\csRad}{#2}%
  \pgfmathsetmacro{\csLead}{1.5625*(#2)}%
  \pgfmathsetmacro{\csLbl}{1.28*(#2)}%
  \begin{scope}[shift={(#1)}]
    \draw (0,0) circle[radius=\csRad cm];
    \draw (0,\csRad) -- (0,\csLead);
    \draw (0,-\csRad) -- (0,-\csLead);
    \node[nlab, anchor=west] at (0,\csLbl)  {#3};
    \node[nlab, anchor=west] at (0,-\csLbl) {#4};
    \coordinate (cs-c) at (0,0);
    \coordinate (cs-a) at (0,\csLead);
    \coordinate (cs-b) at (0,-\csLead);
  \end{scope}%
}
